\subsection{仿射线性映射}

定义 3. 给定附属于同一域 K 上两个向量空间 T、T' 的两个仿射空间 E、E'，E 到 E' 的映射 u 称为仿射线性映射（或仿射映射），如果对 E 的点的每一族 \((x_{t})_{t\in I}\) 以及每一族 \((\lambda_{t})_{t\in I}\) （满足 \(\sum_{i\in I}\lambda_{i}=1\) ），

\[
u \left(\sum_ {i \in I} \lambda_ {i} x _ {i}\right) = \sum_ {i \in I} \lambda_ {i} u (x _ {i}).\tag{3}
\]

命题 6. 设 \(u\) 是 \(E\) 到 \(E'\) 的一个仿射映射。存在一个且仅一个线性映射 \(v\) ，它把 \(T\) 映到 \(T'\) ，使得

\[
u (x + \mathbf {t}) = u (x) + v (\mathbf {t})
\]

对所有 \(x \in \mathbf{E}\) ， \(\mathbf{t} \in \mathbf{T}\) 成立。

设 \(a\) 是 E 中任意一点。映射

\[
\mathbf {t} \mapsto u (a + \mathbf {t}) - u (a)
\]

是 T 到 T' 的一个线性映射 \(v_{a}\) ，因为我们可以写出

\[
\begin{array}{c} a + \lambda \mathbf {t} = \lambda (a + \mathbf {t}) + (1 - \lambda) a \\ a + \mathbf {s} + \mathbf {t} = (a + \mathbf {s}) + (a + \mathbf {t}) - a \end{array}
\]

且由 (3) 可得 \(v_{a}(\lambda \mathbf{t}) = \lambda v_{a}(\mathbf{t})\) 与 \(v_{a}(\mathbf{s} + \mathbf{t}) = v_{a}(\mathbf{s}) + v_{a}(\mathbf{t})\) 。此外，若 \(b\) 是 \(\mathbf{E}\) 中另一点，则 \(v_{a} = v_{b}\) ；因为关系

\[
(a + \mathbf {t}) - a + b = b + \mathbf {t}
\]

蕴含

\[
u (a + \mathbf {t}) - u (a) + u (b) = u (b + \mathbf {t})
\]

也就是 \(u(a + \mathbf{t}) - u(a) = u(b + \mathbf{t}) - u(b)\) 。由此得到 \(v\) 的存在性；唯一性是直接的。

v 称为与 u 相伴的、从 T 到 T' 的线性映射。反之，对 T 到 T' 的每个线性映射 v 以及每对点 \(a \in E\) , \(a' \in E'\) ，立即可验证

\[
x \mapsto a ^ {\prime} + v (x - a)
\]

是从 E 到 \(E'\) 的一个仿射映射，其相伴线性映射为 v。因此，说 u 是 E 到 \(E'\) 的仿射映射，也就意味着：若取 E 中任意一点 a 与点 \(u(a)\) （在 \(E'\) 中）作为原点，则 u 是关于由此所得两个向量空间的一个线性映射。

设 \(\mathbf{E}''\) 是第三个仿射空间， \(\mathbf{T}''\) 是它的平移空间， \(u'\) 是 \(\mathbf{E}'\) 到 \(\mathbf{E}''\) 的一个仿射映射，而 \(v'\) 是一个线性映射（从 \(\mathbf{T}'\) 到 \(\mathbf{T}''\) ），它与 \(u'\) 相伴。显然 \(u' \circ u\) 是 \(\mathbf{E}\) 到 \(\mathbf{E}''\) 的一个仿射映射；此外，对于 \(a \in \mathbf{E}\) 和 \(\mathbf{t} \in \mathbf{T}\) ，

\[
u ^ {\prime} (u (a + \mathbf {t})) = u ^ {\prime} (u (a) + v (\mathbf {t})) = u ^ {\prime} (u (a)) + v ^ {\prime} (v (\mathbf {t}))
\]

因此 \(v'\circ v\) 是 T 到 \(T''\) 的、与 \(u'\circ u\) 相伴的线性映射。仿射映射 u 是双射，必要且充分条件是它的相伴线性映射 v 是双射，且此时 \(u^{-1}\) 是一个仿射映射，其相伴线性映射为 \(v^{-1}\) 。

特别地，E 到自身的仿射双射构成一个群 G，称为 E 的仿射群。将 \(u \in G\) 映射到与 u 相伴的线性映射 v 的映射，由上述可知是 G 到线性群 \(\mathbf{GL}(\mathrm{T})\) 上的同态。若 u 是平移，则 v 是恒等映射，反之亦然。因此，上述同态的核是 E 的平移群 T，从而 T 是 G 的正规子群。

若 \(u \in G\) ，T 的自同构 \(t \mapsto utu^{-1}\) （其中 t 恒同于平移 \(x \mapsto x + t\) ）就是与 u 相伴的线性映射 v。对于 \(x \in E\) 和 \(t \in T\) ，根据定义

\[
x + u \mathbf {t} u ^ {- 1} = u (u ^ {- 1} (x) + \mathbf {t}) = u (u ^ {- 1} (x)) + v (\mathbf {t}) = x + v (\mathbf {t})
\]

因此 \(utu^{-1} = v(t)\) 。

设 \(a \in E\) ， \(G_a\) 表示 \(G\) 中由 \(u \in G\) （满足 \(u(a) = a\) ）组成的子群。若把 \(E\) 与 \(T\) 等同起来（取 \(a\) 为原点），则 \(G_a\) 就与 \(\mathbf{GL}(T)\) 等同。每个 \(u \in G\) 都可以唯一地表示成 \(u = t_1 u_1\) （相应地 \(u = u_2 t_2\) ）的形式，其中 \(u_1, u_2\) 属于 \(G_a\) ， \(t_1, t_2\) 属于 \(T\) ：因为令 \(t_1 = u(a) - a\) ，则 \(u^{-1} t_1 \in G_a\) ，由此得出 \(u_1\) 与 \(t_1\) 的存在性； \(u_2\) 与 \(t_2\) 的存在性同理可得。唯一性由 \(G_a \cap T\) 化为 \(G\) 的单位元这一事实得出。此外，

\[
\mathbf {t} _ {1} u _ {1} = u _ {1} (u _ {1} ^ {- 1} \mathbf {t} _ {1} u _ {1})
\]

由此得 \(u_{2}=u_{1}\) ， \(t_{2}=u_{1}^{-1}t_{1}u_{1}\) 。最后，与 u 和 \(u_{1}\) 相伴的线性映射相同，因此，若如上所述把 \(G_{a}\) 与 \(\mathbf{GL}(T)\) 等同起来，则 \(u_{1}\) 就是从 T 到自身的、与 u 相伴的线性映射。由此可见，G 是 \(G_{a}\) 与 T 的半直积（I，§6，第 1 号）。

设 E，E' 是 K 上的两个仿射空间。在 E（相应地 E'）到 E' 的仿射映射 u 下，E（相应地 E'）的线性簇的直像（相应地逆像）是 E'（相应地 E）的线性簇；u 的秩按定义是 \(u(\mathrm{E})\) 的维数；它等于与 u 相伴的线性映射的秩。若 V，V' 分别是 E，E' 中具有相同有限维数 m 的线性簇，则存在一个仿射映射 \(u\) ，把 \(\mathbf{E}\) 映到 \(\mathbf{E}'\) ，使得 \(u(\mathbf{V}) = \mathbf{V}'\) ：在 \(\mathbf{E}\) 与 \(\mathbf{E}'\) 中分别取 \(\mathbf{V}\) 与 \(\mathbf{V}'\) 的点作为原点，然后在 \(\mathbf{E}\) （相应地 \(\mathbf{E}'\) ）中取一个基，使其前 \(m\) 个向量构成 \(\mathbf{V}\) （相应地 \(\mathbf{V}'\) ）的基，该命题便立即可由 §1 第 11 号命题 17 的推论 3 得出。

由于域 K 典范地具有 K 上的（一维）左向量空间结构，它可以看作一个一维仿射空间。仿射空间 D（在 K 上）到仿射空间 K 的仿射映射也称为仿射线性函数（或仿射函数）。若在 E 中取一点 \(a\) 为原点，则 E 上的每个仿射函数都可以唯一地写成 \(x \mapsto \alpha + v(x)\) ，其中 \(\alpha \in \mathbf{K}\) ，而 \(v\) 是由此所得向量空间 E 上的一个线性形式；因此 E 上的仿射函数构成 K 上一个维数为 \(1 + \dim \mathbf{E}\) 的右向量空间。若 \(u\) 是 E 上一个非常数仿射函数，且 \(\lambda \in \mathbf{K}\) ，则由 \(x \in \mathbf{E}\) （满足方程 \(u(x) = \lambda\) ）组成的集合是一个超平面；反之，对每一个

E 中的超平面 H，存在 E 上的一个仿射函数 \(u_{0}\) ，使得 \(\mathbf{H} = \bar{u}_{0}^{-1}(0)\) ，且每个满足 \(\mathbf{H} = \bar{u}^{-1}(0)\) 的仿射函数 u 都具有形式 \(u_{0}\mu\) ，其中 \(\mu \in K\) （§7，第 5 号，命题 11）。若 u 是 E 上的仿射函数，则方程为 \(u(x) = \alpha\) 与 \(u(x) = \beta\) 的超平面相互平行。

